% This file was converted to LaTeX by Writer2LaTeX ver. 1.9.9
% see http://writer2latex.sourceforge.net for more info
\documentclass{article}
\usepackage{calc,amsmath,amssymb,amsfonts}
\usepackage[T1]{fontenc}
\usepackage[italian]{babel}
\usepackage[style=numeric,backend=biber]{biblatex}
\author{HP}
\date{2026-07-26}
\begin{document}
\section{PROMPT EDITORIALE – LIBRO {\textquotedbl}COMMODORE 64{\textquotedbl}}
Voglio scrivere un libro autobiografico-romanzato sulla nascita della mia passione per l'informatica, utilizzando il
Commodore 64 come filo conduttore. Il libro non deve essere un manuale tecnico né una semplice autobiografia: deve
fondere perfettamente narrativa e didattica, fino al punto che il lettore non percepisca dove finisca il romanzo e dove
inizi l'insegnamento.

L'opera deve ricordare un romanzo di formazione. Il protagonista è un bambino di quinta elementare che, all'inizio degli
anni Ottanta, partecipa a un corso pomeridiano sul Commodore 64 tenuto dal professor Rossi Brunori in una scuola media
di Bari.

Il professor Rossi Brunori non è soltanto un insegnante di informatica. È il mentore del protagonista. Il suo metodo
consiste nel non fornire quasi mai risposte immediate, ma nel guidare gli studenti attraverso domande, esempi concreti
e ragionamenti, affinché siano loro stessi a giungere alle conclusioni. È paziente, ironico, autorevole senza essere
severo e trasmette un autentico amore per la conoscenza.

La parte iniziale del libro racconta il corso di BASIC, ma questo non deve trasformarsi in un manuale completo del
linguaggio. Le lezioni servono esclusivamente ad accendere la curiosità del protagonista e del lettore. Devono essere
affrontati soltanto i concetti indispensabili per comprendere il funzionamento di un programma e per far nascere il
desiderio di programmare.

Gli argomenti iniziali comprendono:

\begin{itemize}
\item accensione del Commodore 64 e schermata READY.; 
\item differenza tra comandi immediati e programma; 
\item numerazione delle linee; 
\item PRINT; 
\item RUN; 
\item GOTO; 
\item il concetto di ciclo infinito; 
\item RUN/STOP come ripresa del controllo del programma. 
\end{itemize}
Dopo questa introduzione il corso deve rapidamente lasciare spazio all'esperienza personale del protagonista.

Il vero apprendimento deve iniziare quando il protagonista ottiene il proprio Commodore 64. Da quel momento ogni nuovo
concetto tecnico nasce da un'esigenza concreta.

L'obiettivo del protagonista non è imparare il BASIC.

L'obiettivo è costruire il proprio videogioco.

Il videogioco rappresenta il motore narrativo dell'intero libro.

Ogni comando BASIC, ogni PEEK, ogni POKE, ogni sprite, ogni interrupt, ogni routine assembler deve comparire soltanto
nel momento in cui diventa necessario per risolvere un problema reale nello sviluppo del gioco.

Il lettore deve avere la sensazione di costruire quel videogioco insieme al protagonista.

Il libro deve accompagnare il lettore dalla prima riga BASIC fino alla realizzazione di un videogioco completo, passando
gradualmente attraverso:

\begin{itemize}
\item grafica; 
\item sprite; 
\item joystick; 
\item collisioni; 
\item animazioni; 
\item sonoro SID; 
\item caricamenti; 
\item gestione della memoria; 
\item ottimizzazioni; 
\item primi approcci all'Assembler 6510. 
\end{itemize}
La progressione tecnica deve risultare completamente naturale.

Ogni nuovo argomento deve nascere perché serve al gioco, mai perché {\textquotedbl}è arrivato il momento di
spiegarlo{\textquotedbl}.

\subsection{Stile narrativo}
La scrittura deve sembrare quella di un romanziere esperto.

Non deve mai ricordare ChatGPT.

Non deve mai ricordare un manuale scolastico.

Il ritmo deve essere fluido, elegante e discorsivo.

Sono vietate:

\begin{itemize}
\item frasi atomiche; 
\item paragrafi costituiti da una o due righe; 
\item elenchi puntati all'interno del manoscritto; 
\item ripetizioni degli stessi concetti; 
\item riassunti inutili; 
\item anticipazioni didascaliche; 
\item frasi tipiche dei testi generati dall'intelligenza artificiale. 
\end{itemize}
I dialoghi devono essere realistici.

Le spiegazioni tecniche devono nascere dai dialoghi e dalle situazioni, non dalla voce narrante.

Quando il narratore adulto interviene, lo fa soltanto per aggiungere una riflessione maturata negli anni, mai per
interrompere la narrazione con lunghe spiegazioni.

\subsection{Filosofia del libro}
Il protagonista non sta imparando un linguaggio di programmazione.

Sta imparando un modo di pensare.

Il Commodore 64 non è il vero protagonista.

Il vero protagonista è la trasformazione di un bambino curioso in uno sviluppatore software.

Il videogioco rappresenta simbolicamente questo percorso di crescita.

Alla fine del libro il lettore dovrà avere la sensazione di aver assistito contemporaneamente:

\begin{itemize}
\item alla nascita di una vocazione; 
\item alla crescita di un programmatore; 
\item alla realizzazione di un videogioco; 
\item a uno spaccato autentico dell'Italia informatica degli anni Ottanta. 
\end{itemize}
L'equilibrio tra narrativa e tecnica dovrà rimanere costante per tutta l'opera.

Ogni pagina dovrà poter essere letta con piacere sia da chi non ha mai visto un Commodore 64, sia da chi lo ha
programmato per anni.


\bigskip
\end{document}
